\documentclass[
	xcolor=dvipsnames,
	10pt, 
	]{beamer}
\graphicspath{{./fig/}}
\RequirePackage{"./sty/UserPackageManager"}

%% TUBS
  \usebackgroundtemplate%
{
	\includegraphics[width=\paperwidth,height=\paperheight]{./bonos/Bonobono}%
}


\title[]{Notes on Baryon Polarization}
\author[\frac{1}{2}uthor]{Bono}
\institute{Bono University}
\date{\today}
\subtitle{\today}


\begin{document}


%%%%%%%%%%%%%%% Title page 
\begin{frame}[t,plain] % Cover slide
       \titlepage
\end{frame}
\begin{frame}{Spin State Representation of Baryon Decay}
	\begin{block}{}
		Consider a decay process where spin-$\frac{1}{2}$ baryon($\Xi^-$)  decays into another spin-$\frac{1}{2}$ baryon ($\Lambda$) and a spin-0 meson($\pi^-$). In quantum mechanics, this decay $\Xi^-\to \Lambda\pi^-$  can be modeled as a transition from the initial state $\ket{\Xi^-}$ to the final state $\ket{\Lambda,\pi^-}$. Since $\pi^-$ is a spin-0 meson, we focues on the orbital angular momentum carried by $\pi^-$. In the S-wave transition  $\pi^-$ carries no angular momentum($l=0$), while in the P-wave transition it carries an angular momentum of $l=1$.  The initial state of $\Xi^-$ can be defined as		$\Xi^-=\ket{\frac{1}{2},\frac{1}{2}}$. The decay process is then represented as:
		\begin{align}
			\ket{\frac{1}{2},\frac{1}{2}}_\Xi &= A \ket{\frac{1}{2},\frac{1}{2}}_\Lambda\ket{0,0}_\pi Y_0^0 +B[C_{\frac{1}{2},\frac{1}{2},1,0}^{\frac{1}{2},\frac{1}{2}}\ket{\frac{1}{2},\frac{1}{2}}_\Lambda\ket{1,0}\pi Y_1^0\nonumber\\
			&+C_{\frac{1}{2},-\frac{1}{2},1,1}^{\frac{1}{2},\frac{1}{2}}\ket{\frac{1}{2},-\frac{1}{2}}_\Lambda\ket{1,1}\pi Y_1^1]\label{Decay}
		\end{align}
	\end{block}
\end{frame}
\begin{frame}{Spin State Representation of Baryon Decay}
	\begin{block}{}
	Here A and B represent the transition amplitudes of S-wave and P-wave, respectively. The terms $C_{j1,m1,j2,m2}^{J,M}$  are Clebsch-Gordan coefficients, and $Y_l^m(\theta,\phi)$s denote spherical harmonics. Note that A and B are complex numbers. Focusing on the spin state of $\Lambda$, we rewrite \eqref{Decay} as:
	\begin{align}
		\ket\Lambda = (a + b\cos\theta)\ket{\frac{1}{2},\frac{1}{2}}_\Lambda + be^{i\phi}\sin\theta \ket{\frac{1}{2},-\frac{1}{2}}_\Lambda.\label{State}
	\end{align}
	The decay spectrum of $\Lambda$ can be expressed as:
	\begin{align}
		\bra\Lambda\ket\Lambda = (a+b\cos\theta)^* (a+b\cos\theta)\Lambda_{\uparrow}+ (be^{i\phi})^*be^{i\phi}\Lambda_{\downarrow}\label{Lambda}
	\end{align}
	Here, $\Lambda_{\uparrow}$ and $\Lambda_{\downarrow}$ represent the spin-up and spin-down states of $\Lambda$, respectively.
\end{block}
\end{frame}
\begin{frame}{Baryon Decay Parameters}
	\begin{block}{}
		If we not distinguish $\Lambda_{\uparrow}$ and $\Lambda_{\downarrow}$, \eqref{Lambda} turns into:
		\begin{align}
			\Lambda(\theta,\phi)=\abs{a}^2+\abs{b}^2+2\re[a^*b]\cos\theta = 1+\alpha\cos\theta.
		\end{align}
		The wavefunction should be normalized at the final calculation, but at the moment we will require $\abs{a}^2+\abs{b}^2=1$ and introduce decay parameters:
		\begin{align}
			\alpha= \re[a^*b], \beta = \im[a^*b],\gamma = \abs{a}^2-\abs{b}^2.
		\end{align}
		Remind that our coordinate system here is the Center-of-Mass frame of $\Xi^-$, with spin pointing the z axis as we already set the initial state as $\ket\Xi=\ket{\frac{1}{2},\frac{1}{2}}$. $\theta$ and $\phi$ corresponds to the spherical variables of $\Lambda$ momentum at CoM of $\Xi^-$.
	\end{block}
\end{frame}

\begin{frame}{Spin State of Daughter Baryon}
	\begin{block}{}
		Now time has come to calculate the spin state of $\Lambda$ in \textbf{CoM of }$\mathbf{\Xi^-}$. The spinnor representation of $\Lambda$, as given in \eqref{State}, is:
		\begin{align}
			\ket\Lambda = \begin{pmatrix}
		 a + b\cos\theta\\
		  be^{i\phi}\sin\theta 
			\end{pmatrix}\label{Spinnor}
		\end{align}
		and the direction of spin, or, the \textit{polarization} of $\Lambda$ is
		\begin{align}
			\vec P=\bra\Lambda \hat S\ket \Lambda.
		\end{align}
		Here we defined the spin-operator $\hat S = \{\sigma_x,\sigma_y,\sigma_z\}$, where $\sigma_i$s are the \textit{Pauli matrices}.
		Our task is to calculate each component of $\vec P$, \textit{at Center of Mass of $\Lambda$}.
	\end{block}
\end{frame}
\begin{frame}{Calculating $P_x$}
	\begin{block}{}
		We start with calculating $P_x$. $\sigma_x = \dmat 0 1 1 0$ and
		\begin{align}
			P_x& =\bra\Lambda \dmat 0 1 1 0 \ket \Lambda\nonumber\\
			& = (a +b \cos\theta)^*be^{i\phi}\sin\theta+(be^{i\phi}\sin\theta)^*(a +b \cos\theta)\nonumber\\
			&= (a^*be^{i\phi}+ab^*e^{-i\phi})\sin\theta + \abs{b}^2\sin\theta\cos\theta(e^{i\phi}+e^{-i\phi})\nonumber\\
			&=(2\re[a^*b]\cos\phi-2 \im[a^*b]\sin\phi )\sin\theta+2\abs{b}^2\sin\theta\cos\theta\cos\phi\nonumber\\
			&= [\alpha\cos\phi-\beta\sin\phi+ (1-\gamma)\cos\theta\cos\phi]\sin\theta.
		\end{align}
		Here we used $2\abs{b}^2 = \abs{a}^2+\abs{b}^2 - (\abs{a}^2-\abs{b}^2)=1-\gamma$.
	\end{block}
\end{frame}
\begin{frame}{Calculating $P_y$}
	\begin{block}{}
		\begin{align}
	P_y& =\bra\Lambda \dmat {0}{-i}{i}{0} \ket \Lambda\nonumber\\
	& =-i (a +b \cos\theta)^*be^{i\phi}\sin\theta+i(be^{i\phi}\sin\theta)^*(a +b \cos\theta)\nonumber\\
	&= -i[(a^*be^{i\phi}-ab^*e^{-i\phi})\sin\theta + \abs{b}^2\sin\theta\cos\theta(e^{i\phi}-e^{-i\phi})]\nonumber\\
	&=-i[(2i\im[a^*b]\cos\phi+2 i\re[a^*b]\sin\phi )\sin\theta+2i\abs{b}^2\sin\theta\cos\theta\sin\phi]\nonumber\\
	&= [\beta\cos\phi+\alpha\sin\phi+ (1-\gamma)\cos\theta\sin\phi]\sin\theta.
\end{align}
	\end{block}
\end{frame}
\begin{frame}{Calculating $P_z$}
	\begin{block}{}
		\begin{align}
			P_z& =\bra\Lambda \dmat {1}{0}{0}{-1} \ket \Lambda\nonumber\\
			&= (a+b\cos\theta)^*(a+b\cos\theta)-(b^*e^{-i\phi})be^{i\phi}\nonumber\\
			&= \abs {a}^2 + (a^*b+ab^*)\cos\theta + \abs{b}^2\cos^2\theta -\abs{b}^2\sin^2\theta\nonumber\\
			&= \frac{1}{2}(1+\gamma) + 2\re[a^*b]\cos\theta  + \frac{\cos^2\theta-\sin^2\theta}{2}(1-\gamma)\nonumber\\
			&=\frac{1}{2}(1+\gamma) + \alpha\cos\theta  + \frac{\cos^2\theta-\sin^2\theta}{2}(1-\gamma)
.		\end{align}
	\end{block}
\end{frame}
\end{document}
